\documentclass[aspectratio=169,11pt]{beamer}

\usetheme{Madrid}
\usecolortheme{default}

\usepackage[utf8]{inputenc}
\usepackage[T1]{fontenc}
\usepackage[brazil]{babel}
\usepackage{lmodern}
\usepackage{amsmath}
\usepackage{booktabs}
\usepackage{array}
\usepackage{multirow}
\usepackage{tikz}
\usepackage{pgfplots}
\usepackage{pgf-pie}
\usetikzlibrary{positioning,arrows.meta,shapes.geometric}
\pgfplotsset{compat=1.18}

\definecolor{seplanblue}{RGB}{20,52,100}
\definecolor{seplanblue2}{RGB}{32,78,145}
\definecolor{seplangold}{RGB}{200,145,20}
\definecolor{seplangreen}{RGB}{34,130,64}
\definecolor{seplanred}{RGB}{180,40,40}
\definecolor{seplanlight}{RGB}{236,242,248}

\setbeamercolor{palette primary}{bg=seplanblue,fg=white}
\setbeamercolor{palette secondary}{bg=seplanblue2,fg=white}
\setbeamercolor{palette tertiary}{bg=seplangold,fg=white}
\setbeamercolor{frametitle}{bg=seplanblue,fg=white}
\setbeamercolor{title}{fg=seplanblue}
\setbeamercolor{structure}{fg=seplanblue}
\setbeamercolor{block title}{bg=seplanblue2,fg=white}
\setbeamercolor{block body}{bg=seplanlight,fg=black}
\setbeamercolor{itemize item}{fg=seplanblue}
\setbeamercolor{itemize subitem}{fg=seplangold}

\setbeamertemplate{navigation symbols}{}
\setbeamertemplate{footline}{%
  \leavevmode%
  \hbox{%
  \begin{beamercolorbox}[wd=.78\paperwidth,ht=2.5ex,dp=1.2ex,leftskip=1em]{palette primary}%
    \insertshorttitle
  \end{beamercolorbox}%
  \begin{beamercolorbox}[wd=.22\paperwidth,ht=2.5ex,dp=1.2ex,rightskip=1em,right]{palette secondary}%
    \insertframenumber/\inserttotalframenumber
  \end{beamercolorbox}}%
}

\title[IAET-RR]{IAET-RR: defesa técnica do projeto}
\subtitle{Indicador de Atividade Econômica Trimestral de Roraima}
\author{Yuri Cesar de Lima e Silva}
\institute{SEPLAN/RR -- CGEES/DIEAS}
\date{Abril de 2026}

\begin{document}

\begin{frame}[plain]
  \titlepage
  \vspace{-0.4cm}
  \begin{center}
    \small Apresentação técnica para discussão metodológica
  \end{center}
\end{frame}

\begin{frame}{Qual lacuna o projeto resolve}
\begin{itemize}
  \item O IBGE divulga as Contas Regionais em frequência anual e com defasagem.
  \item A gestão precisa de leitura trimestral para monitoramento, comunicação e decisão.
  \item O projeto cria uma ponte entre dados administrativos de alta frequência e o arcabouço anual das Contas Regionais.
\end{itemize}

\vspace{0.3cm}
\begin{block}{Mensagem central}
O projeto não ``inventa um PIB trimestral''. Ele trimestraliza informação econômica observável, disciplinada por benchmark anual oficial.
\end{block}
\end{frame}

\begin{frame}{O produto estatístico}
\begin{columns}[T]
\column{0.48\textwidth}
\textbf{Saídas principais}
\begin{itemize}
  \item índice trimestral de atividade econômica em volume
  \item série dessazonalizada
  \item VAB nominal trimestral
  \item PIB nominal trimestral
  \item PIB real trimestral
  \item VAB nominal setorial
  \item dashboard interativo
\end{itemize}

\column{0.48\textwidth}
\textbf{Lógica de uso}
\begin{itemize}
  \item conjuntura de curtíssimo prazo
  \item monitoramento setorial
  \item comunicação institucional
  \item base para boletins e notas técnicas
\end{itemize}
\end{columns}
\end{frame}

\begin{frame}{Princípio metodológico}
\begin{enumerate}
  \item construir indicadores trimestrais por setor;
  \item agregar com pesos compatíveis com o ano-base;
  \item impor coerência com benchmark anual oficial;
  \item derivar séries nominais e reais consistentes.
\end{enumerate}

\vspace{0.3cm}
\begin{block}{Intuição}
As Contas Regionais dizem qual é o nível anual correto. As proxies trimestrais dizem como esse nível se distribui dentro do ano.
\end{block}
\end{frame}

\begin{frame}{Estrutura setorial do projeto}
\begin{columns}[T]
\column{0.52\textwidth}
\small
\begin{tabular}{lr}
\toprule
\textbf{Setor} & \textbf{Peso 2020} \\
\midrule
Adm. Pública & 45,01\% \\
Serviços & 36,46\% \\
Indústria & 11,63\% \\
Agropecuária & 6,89\% \\
\bottomrule
\end{tabular}

\vspace{0.35cm}
\begin{itemize}
  \item A elevada participação de Adm. Pública exige modelagem aderente à estrutura de Roraima.
  \item Isso diferencia o projeto de leituras genéricas baseadas apenas em proxies de mercado.
\end{itemize}

\column{0.44\textwidth}
\centering
\begin{tikzpicture}
\pie[
  text=legend,
  radius=1.75,
  color={seplanblue,seplanblue2,seplangold,seplangreen}
]{45.01/Adm.\ Pública,36.46/Serviços,11.63/Indústria,6.89/Agropecuária}
\end{tikzpicture}
\end{columns}
\end{frame}

\begin{frame}{Agregação em volume}
\[
I_t = \sum_{i=1}^{n} w_i^{(2020)} \cdot I_{i,t}
\]

onde:
\begin{itemize}
  \item \(I_t\) = índice agregado no trimestre \(t\)
  \item \(I_{i,t}\) = índice do setor \(i\) no trimestre \(t\)
  \item \(w_i^{(2020)}\) = peso do setor \(i\) no VAB nominal do ano-base 2020
\end{itemize}

\begin{block}{Intuição}
A estrutura de ponderação fica fixa no ano-base; o que varia trimestre a trimestre é o desempenho de cada setor.
\end{block}
\end{frame}

\begin{frame}{O papel do Denton-Cholette}
\textbf{Problema}
\begin{itemize}
  \item temos benchmark anual oficial;
  \item precisamos de trajetória trimestral consistente com esse benchmark.
\end{itemize}

\vspace{0.2cm}
\textbf{Condição de coerência}

Para séries em índice:
\[
\frac{1}{4}\sum_{q=1}^{4} x_{y,q} = B_y
\]

Para séries de fluxo:
\[
\sum_{q=1}^{4} x_{y,q} = B_y
\]

\vspace{0.2cm}
O método encontra \(x_t\) preservando o máximo possível a dinâmica do indicador trimestral \(z_t\).
\end{frame}

\begin{frame}{Denton-Cholette: intuição econômica}
\begin{itemize}
  \item o benchmark anual fixa o envelope da série;
  \item a proxy trimestral informa o desenho intra-anual.
\end{itemize}

\begin{block}{Em termos práticos}
\begin{itemize}
  \item sem benchmark anual, a proxy pode carregar viés de nível;
  \item sem proxy trimestral, a trimestralização tende a ser artificial;
  \item com Denton, a série trimestral respeita o dado anual e preserva o perfil observado.
\end{itemize}
\end{block}
\end{frame}

\begin{frame}{Por que isso é importante}
\begin{itemize}
  \item evita que o indicador trimestral descole das Contas Regionais;
  \item reduz arbitrariedade na distribuição intra-anual;
  \item produz uma série defensável tecnicamente.
\end{itemize}

\begin{alertblock}{Mensagem para gestão}
O projeto antecipa a leitura da economia sem abrir mão da disciplina estatística do benchmark oficial.
\end{alertblock}
\end{frame}

\begin{frame}{Construção do VAB nominal}
\[
I_t^{nom} = I_t^{real} \times \frac{D_t}{100}
\]

onde:
\begin{itemize}
  \item \(I_t^{nom}\) = índice nominal
  \item \(I_t^{real}\) = índice real
  \item \(D_t\) = deflator implícito trimestral
\end{itemize}

\textbf{Como o deflator é obtido}
\begin{enumerate}
  \item calcula-se o deflator anual a partir das Contas Regionais;
  \item trimestraliza-se o deflator anual com Denton-Cholette;
  \item usa-se o IPCA como indicador auxiliar intra-anual.
\end{enumerate}
\end{frame}

\begin{frame}{Construção do PIB nominal}
\[
PIB = VAB + ILP
\]

onde:
\begin{itemize}
  \item \(VAB\) = valor adicionado bruto
  \item \(ILP\) = impostos líquidos sobre produtos
\end{itemize}

\begin{itemize}
  \item O VAB nominal trimestral é estimado a partir do bloco real e do deflator.
  \item O ILP trimestral é distribuído com benchmark anual e proxy tributária.
  \item O PIB nominal trimestral resulta da soma dos dois componentes.
\end{itemize}
\end{frame}

\begin{frame}{Construção do PIB real}
\textbf{Etapa 1: série preliminar}
\[
PIB_t^{real,prel} = \frac{PIB_t^{nom}}{D_t/100}
\]

\textbf{Etapa 2}
\begin{itemize}
  \item calcular o índice anual preliminar;
  \item comparar com o benchmark oficial de crescimento real do PIB.
\end{itemize}

\textbf{Etapa 3}
\begin{itemize}
  \item ancorar a trajetória trimestral por Denton-Cholette;
  \item impor coincidência da média anual com o benchmark oficial do PIB real.
\end{itemize}

\begin{block}{Ponto metodológico}
Isso reduz o problema de a variação real dos impostos sobre produtos contaminar a leitura do PIB real.
\end{block}
\end{frame}

\begin{frame}{Benchmarking do PIB real}
\centering
\begin{tabular}{lrr}
\toprule
\textbf{Ano} & \textbf{Projeto} & \textbf{CR IBGE} \\
\midrule
2021 & 8,40\% & 8,40\% \\
2022 & 11,30\% & 11,30\% \\
2023 & 4,20\% & 4,20\% \\
\bottomrule
\end{tabular}

\vspace{0.4cm}
\begin{block}{Leitura}
Nos anos com benchmark oficial disponível, o crescimento real anual do PIB do projeto coincide com as Contas Regionais.
\end{block}
\end{frame}

\begin{frame}{Resultados: PIB nominal anual}
\begin{columns}[T]
\column{0.40\textwidth}
\small
\begin{tabular}{lr}
\toprule
\textbf{Ano} & \textbf{R\$ bi} \\
\midrule
2020 & 16,02 \\
2021 & 18,20 \\
2022 & 21,10 \\
2023 & 25,12 \\
2024 & 28,86 \\
2025 & 32,59 \\
\bottomrule
\end{tabular}

\column{0.58\textwidth}
\begin{tikzpicture}
\begin{axis}[
  ybar,
  bar width=12pt,
  ymin=0,ymax=35,
  width=\textwidth,
  height=5.8cm,
  ylabel={R\$ bi},
  symbolic x coords={2020,2021,2022,2023,2024,2025},
  xtick=data,
  nodes near coords,
  every node near coord/.append style={font=\scriptsize, rotate=90, anchor=west},
  nodes near coords align={vertical},
  enlarge x limits=0.08,
  axis x line*=bottom,
  axis y line*=left,
  ymajorgrids=true,
  grid style={dashed,gray!30}
]
\addplot[fill=seplanblue] coordinates {
  (2020,16.02) (2021,18.20) (2022,21.10) (2023,25.12) (2024,28.86) (2025,32.59)
};
\end{axis}
\end{tikzpicture}
\end{columns}
\end{frame}

\begin{frame}{Resultados: crescimento real do PIB}
\begin{columns}[T]
\column{0.38\textwidth}
\small
\begin{tabular}{lr}
\toprule
\textbf{Ano} & \textbf{\%} \\
\midrule
2021 & 8,40 \\
2022 & 11,30 \\
2023 & 4,20 \\
2024 & 6,54 \\
2025 & 7,54 \\
\bottomrule
\end{tabular}

\column{0.60\textwidth}
\begin{tikzpicture}
\begin{axis}[
  ybar,
  bar width=14pt,
  ymin=0,ymax=12,
  width=\textwidth,
  height=5.8cm,
  ylabel={\%},
  symbolic x coords={2021,2022,2023,2024,2025},
  xtick=data,
  nodes near coords,
  every node near coord/.append style={font=\scriptsize},
  enlarge x limits=0.10,
  axis x line*=bottom,
  axis y line*=left,
  ymajorgrids=true,
  grid style={dashed,gray!30}
]
\addplot[fill=seplangold] coordinates {
  (2021,8.40) (2022,11.30) (2023,4.20) (2024,6.54) (2025,7.54)
};
\end{axis}
\end{tikzpicture}
\end{columns}
\end{frame}

\begin{frame}{Resultados: estrutura setorial estimada de 2025}
\begin{columns}[T]
\column{0.52\textwidth}
\small
\begin{tabular}{lrr}
\toprule
\textbf{Setor} & \textbf{R\$ bi} & \textbf{\%} \\
\midrule
Adm. Pública & 15,74 & 47,4 \\
Serviços & 9,08 & 27,4 \\
Indústria & 4,97 & 15,0 \\
Agropecuária & 3,39 & 10,2 \\
\bottomrule
\end{tabular}

\column{0.42\textwidth}
\centering
\begin{tikzpicture}
\pie[
  text=legend,
  radius=1.7,
  color={seplanblue,seplanblue2,seplangold,seplangreen}
]{47.4/Adm.\ Pública,27.4/Serviços,15.0/Indústria,10.2/Agropecuária}
\end{tikzpicture}
\end{columns}
\end{frame}

\begin{frame}{Arquitetura do pipeline}
\begin{enumerate}
  \item construir setores reais;
  \item agregar o índice geral;
  \item dessazonalizar;
  \item validar;
  \item gerar VAB nominal;
  \item gerar PIB nominal;
  \item gerar PIB real;
  \item gerar VAB nominal setorial;
  \item publicar no painel.
\end{enumerate}

\begin{block}{Ganho institucional}
Processo replicável, auditável e com baixo custo marginal de atualização.
\end{block}
\end{frame}

\begin{frame}{O dashboard}
\begin{itemize}
  \item leitura do índice geral;
  \item abertura por setor;
  \item PIB nominal;
  \item PIB real;
  \item VAB nominal setorial;
  \item exportação de bases.
\end{itemize}

\vspace{0.3cm}
\begin{block}{Valor}
O painel reduz o atrito entre produção técnica e uso gerencial, transformando o pipeline estatístico em produto de comunicação pública.
\end{block}
\end{frame}

\begin{frame}{Próximas entregas}
\textbf{Curto prazo}
\begin{itemize}
  \item boletins trimestrais de conjuntura;
  \item notas executivas para alta gestão;
  \item consolidação narrativa do dashboard.
\end{itemize}

\textbf{Médio prazo}
\begin{itemize}
  \item relatórios temáticos por setor;
  \item comparações com CR, IBCR, IBC-Br e arrecadação;
  \item rotina institucionalizada de atualização.
\end{itemize}
\end{frame}

\begin{frame}{Como vender o projeto institucionalmente}
\begin{columns}[T]
\column{0.48\textwidth}
\begin{block}{Argumento técnico}
\begin{itemize}
  \item aderência metodológica às Contas Regionais;
  \item benchmark anual oficial;
  \item separação entre séries reais e nominais;
  \item solução explícita para ancoragem do PIB real.
\end{itemize}
\end{block}

\column{0.48\textwidth}
\begin{block}{Argumento de gestão}
\begin{itemize}
  \item ganho de tempestividade;
  \item ganho de inteligibilidade;
  \item ganho de autonomia analítica do estado.
\end{itemize}
\end{block}
\end{columns}
\end{frame}

\begin{frame}{Mensagem final}
\begin{itemize}
  \item o IBGE fornece o benchmark anual;
  \item o projeto fornece a leitura trimestral;
  \item o dashboard fornece usabilidade;
  \item os relatórios futuros fornecem narrativa econômica oficial.
\end{itemize}

\vspace{0.4cm}
\begin{alertblock}{Proposta de valor}
Colocar a SEPLAN/RR em posição de referência na leitura conjuntural da economia estadual.
\end{alertblock}
\end{frame}

\end{document}
